\documentclass[a4paper]{article}

\usepackage[spanish]{babel}
\usepackage{graphicx}
\usepackage{amsmath, amssymb, mathrsfs}
\usepackage[margin=2cm]{geometry}
\usepackage{fancyhdr}
\usepackage{enumerate}
\usepackage[shortlabels]{enumitem}
\usepackage{parskip}
\usepackage[most]{tcolorbox}
\usepackage[hidelinks]{hyperref}
\usepackage{float}
\usepackage{booktabs}
\usepackage{mathrsfs}
\usepackage{tikz}
\usetikzlibrary{arrows.meta, calc, decorations.pathreplacing}
\usepackage{pgfplots}
\pgfplotsset{compat=1.18}


% cabecera
\pagestyle{fancy}
\fancyhead[l]{Física Matemática}
\fancyhead[c]{Parcial 3}
\fancyhead[r]{\today}
\fancyfoot[c]{\thepage}
\renewcommand{\headrulewidth}{0.2pt} % linea horizontal


\begin{document}


\section{Problema 3}
Considere un cascarón esférico de radio $R$ cuyo potencial electrostático sobre la superficie es 
$$\phi(R,\theta) = V_0 \sin^2\left(\frac{\theta}{2}\right)$$
\begin{enumerate}[a)]
\item Determinar el potencial eléctrico $\phi(r,\theta)$ en el interior de la esfera.
\item Hallar las componentes radial $E_r$ y transversal $E_\theta$ del campo eléctrico en el interior de la esfera
\end{enumerate}

\subsection{Solucion Problema 3}
\textbf{item a)}

Conocemos la solución de $\phi$ en esféricas de manera general. Para este problema, claramente debemos imponer que $B_l = 0$, entonces nos queda:

$$
\phi(r,\theta) = \sum_{l=0}^{\infty} A_l r^l P_l(\cos\theta)
$$

Vamos a desarrollar la solución a este problema de manera inteligente. Sabemos que en alguna parte de la solución del problema debemos recurrir a la condición de continuidad para $r<R$ con $r=R$ y sabemos que para un problema con simetría esférica, los primeros términos de una expansión con Legendre son $1$ y $\cos\theta$. Entonces miramos la forma de $V\Bigg|_{r=R}$ y notamos que es un $\sin^2\left(\frac{\theta}{2}\right)$. Entonces recordamos que $\sin^2(x)$ se puede llevar a una forma con $\cos(2x)$, por lo tanto, $\sin^2\left(\frac{\theta}{2}\right)$ se puede llevar a una forma con $\cos\theta$, lo cual nos es muy útil en una expansión de Legendre, porque tenemos términos de $\cos\theta$. Teniendo en cuenta este razonamiento, entonces miremos la expansión de $\phi(R,\theta)$ con legendre y comparemos los coeficientes.

Primero expandimos la expresión del potencial en $R$.

$$\phi(R,\theta) = \sum_{l=0}^{\infty} A'_l R^l P_l(\cos\theta) = \frac{V_0}{2} - \frac{V_0}{2R}\cos\theta$$

Ahora por condición de continuidad, necesitamos que 

$$\phi(r,\theta) = \sum_{l=0}^{\infty} A_l r^l P_l(\cos\theta) = \sum_{l=0}^{\infty} A'_l R^l P_l(\cos\theta) = \phi(R,\theta)$$

Con esto obtenemos los primeros términos:

\begin{itemize}
	\item $A_0 = \frac{V_0}{2}$
	\item $A_1 = - \frac{V_0}{2R}$
	\item $A_l = 0\qquad \forall l \geq 2$
\end{itemize}

 Lo aplicamos expandiendo los primeros dos términos.

$$\phi(r,\theta) = \sum_{l=0}^{\infty} A_l r^l P_l(\cos\theta) = \frac{V_0}{2} - \frac{V_0}{2R} r \cos\theta + \sum_{l=2}^{\infty} A_l r^l P_l(\cos\theta) = \phi(R,\theta)$$


Y quedamos con 

$$\phi(r,\theta) = \frac{V_0}{2} - \frac{V_0}{2R}r\cos\theta$$

\textbf{item b)}

Ahora veamos la expresión para el campo eléctrico. Sabemos que 

$$\vec{E} = -\vec{\nabla}_\phi$$ 

y recordemos que en este sistema de coordenadas, 

$$\vec{\nabla} = \frac{\partial}{\partial r} \hat{e_r} + \frac{1}{r}\frac{\partial}{\partial \theta} \hat{e_\theta}$$

Con esto, calculamos el campo eléctrico y nos queda:

$$\vec{E} = \frac{V_0}{2R}\left(\cos\theta \hat{e_r} - \sin\theta \hat{e_\theta}\right) = \frac{V_0}{2R}\hat{e_k}$$


\subsubsection*{item c)}

Para la región exterior ($r > R$), claramente debemos imponer que $A_{l}=0$ para que el potencial sea finito en el infinito ($r \to \infty$), entonces nos queda:
\begin{equation}
\phi(r,\theta) = \sum_{l=0}^{\infty} \frac{B_l}{r^{l+1}} P_l(\cos\theta)
\end{equation}


Recordamos que los primeros polinomios de Legendre son $P_0(\cos\theta) = 1$ y $P_1(\cos\theta) = \cos\theta$. Teniendo en cuenta este razonamiento, miremos la expansión de $\phi(R,\theta)$ con Legendre y comparemos los coeficientes en $r = R$:
\begin{equation}
\phi(R,\theta) = \sum_{l=0}^{\infty} \frac{B_l}{R^{l+1}} P_l(\cos\theta) = \frac{V_{0}}{2}P_0(\cos\theta) - \frac{V_{0}}{2}P_1(\cos\theta)
\end{equation}

Al comparar término a término, obtenemos los coeficientes correspondientes:
\begin{align}
\frac{B_0}{R^1} = \frac{V_0}{2} &\implies B_0 = \frac{V_0 R}{2} \\
\frac{B_1}{R^2} = -\frac{V_0}{2} &\implies B_1 = -\frac{V_0 R^2}{2} \\
B_l = 0 &\quad \forall l \ge 2
\end{align}

Sustituyendo estos coeficientes en la expansión original para el exterior, nos queda:
\begin{equation}
\phi(r,\theta) = \frac{V_{0}R}{2r} - \frac{V_{0}R^{2}}{2r^{2}}\cos\theta
\end{equation}

\subsubsection*{item d)}

Ahora veamos la expresión para el campo eléctrico. Sabemos que
\begin{equation}
\vec{E} = -\vec{\nabla}\phi
\end{equation}
y recordemos que en este sistema de coordenadas, el operador gradiente para componentes con simetría azimutal es:
\begin{equation}
\vec{\nabla} = \frac{\partial}{\partial r}\hat{e}_{r} + \frac{1}{r}\frac{\partial}{\partial\theta}\hat{e}_{\theta}
\end{equation}

Calculamos la componente radial ($E_r$):
\begin{equation}
E_r = -\frac{\partial \phi}{\partial r} = -\left( -\frac{V_0 R}{2r^2} + \frac{V_0 R^2}{r^3}\cos\theta \right) = \frac{V_0 R}{2r^2} - \frac{V_0 R^2}{r^3}\cos\theta
\end{equation}

Calculamos la componente transversal ($E_\theta$):
\begin{equation}
E_\theta = -\frac{1}{r}\frac{\partial \phi}{\partial \theta} = -\frac{1}{r}\left( \frac{V_0 R^2}{2r^2}\sin\theta \right) = -\frac{V_0 R^2}{2r^3}\sin\theta
\end{equation}

Con esto, el campo eléctrico en el exterior ($r > R$) queda determinado por:
\begin{equation}
\vec{E} = \left( \frac{V_0 R}{2r^2} - \frac{V_0 R^2}{r^3}\cos\theta \right)\hat{e}_{r} - \left( \frac{V_0 R^2}{2r^3}\sin\theta \right)\hat{e}_{\theta}
\end{equation}
\bibliographystyle{plainurl}
\bibliography{bibliografia} 
\end{document}
